\documentclass[12pt,a4paper]{report}

\usepackage{fontspec}
\usepackage{polyglossia}
\setdefaultlanguage{romanian}
\usepackage[margin=2.25cm]{geometry}
\usepackage{setspace}
\usepackage{enumitem}
\usepackage{longtable}
\usepackage{array}
\usepackage{booktabs}
\usepackage{xcolor}
\usepackage{listings}
\usepackage[
    colorlinks,
    linkcolor={black},
    citecolor={black},
    urlcolor={blue}
]{hyperref}

\onehalfspacing
\setlength{\emergencystretch}{3em}
\setlength{\parskip}{0.15em}
\setlist{nosep,leftmargin=*}

\definecolor{codebackground}{rgb}{0.96,0.96,0.96}
\definecolor{codecomment}{rgb}{0.23,0.45,0.30}
\definecolor{codekeyword}{rgb}{0.10,0.20,0.55}
\definecolor{codestring}{rgb}{0.55,0.15,0.15}

\lstdefinestyle{codestyle}{
    backgroundcolor=\color{codebackground},
    basicstyle=\ttfamily\scriptsize,
    breaklines=true,
    frame=single,
    columns=fullflexible,
    keepspaces=true,
    showstringspaces=false,
    numbers=left,
    numberstyle=\tiny\color{gray},
    keywordstyle=\color{codekeyword}\bfseries,
    commentstyle=\color{codecomment},
    stringstyle=\color{codestring},
    tabsize=4
}
\lstset{style=codestyle}

\newcommand{\filetitle}[1]{\section*{\texttt{\detokenize{#1}}}}
\newcommand{\chunk}[3]{\subsection*{Liniile #2--#3}\lstinputlisting[language=Python,firstline=#2,lastline=#3,firstnumber=#2]{#1}}
\newcommand{\cmdchunk}[3]{\subsection*{Liniile #2--#3}\lstinputlisting[language=bash,firstline=#2,lastline=#3,firstnumber=#2]{#1}}

\title{Explicarea fisierelor responsabile de interfata real-time}
\author{}
\date{}

\begin{document}

\maketitle
\tableofcontents
\clearpage

\chapter{Privire de ansamblu}

Interfata real-time este formata dintr-un grup mic de fisiere care colaboreaza
pentru a porni aplicatia, a captura semnalul audio, a-l procesa, a-l reda si a-l
afisa vizual. In continuare sunt explicate fisierele principale, in ordinea in
care sunt intalnite in fluxul de executie:

\begin{longtable}{p{0.37\textwidth}p{0.55\textwidth}}
\toprule
\textbf{Fisier} & \textbf{Rol in interfata real-time} \\
\midrule
\endhead
\path{launch_realtime_demo.cmd} &
Script Windows de lansare rapida. Alege checkpoint-ul, pregateste directoarele
de recording/snapshot si porneste aplicatia Python. \\
\path{src/tools/realtime_demo.py} &
Entrypoint-ul Python. Parseaza argumentele CLI, construieste enhancer-ul,
interfata si motorul audio. \\
\path{src/tools/realtime/__init__.py} &
Reexporta componentele publice ale pachetului \texttt{realtime}, ca importurile
din entrypoint sa ramana curate. \\
\path{src/tools/realtime/common.py} &
Contine utilitare comune: normalizare audio, statistici, recording, sursa
sintetica \texttt{speech+noise} si selectia dispozitivelor audio. \\
\path{src/tools/realtime/engine.py} &
Motorul efectiv: gestioneaza callback-urile audio, cozile, thread-ul worker si
aplicarea metodelor de denoising. \\
\path{src/tools/realtime/ui.py} &
Interfata Tkinter/Matplotlib: controale, grafice, spectrograme, VAD live si
modelul de cereri \texttt{pending-consume}. \\
\bottomrule
\end{longtable}

Explicatiile sunt organizate pe intervale continue de linii. Pentru liniile
goale sau pentru blocurile de import similare, comentariul este grupat; pentru
liniile care schimba comportamentul aplicatiei, explicatia este detaliata.

\chapter{Launcher-ul Windows}

\filetitle{launch_realtime_demo.cmd}

\cmdchunk{../launch_realtime_demo.cmd}{1}{8}

\begin{description}
    \item[Liniile 1--2.] Dezactiveaza afisarea comenzilor in consola si porneste
    un context local de variabile, astfel incat setarile scriptului sa nu ramana
    in mediul sistemului dupa terminare.
    \item[Liniile 4--8.] Scriptul se muta in directorul in care se afla el
    insusi, apoi defineste radacina proiectului, checkpoint-ul folosit de
    MaskNet, directoarele pentru recording si snapshot-uri si un director-lock
    in \texttt{\%TEMP\%}. Lock-ul impiedica pornirea mai multor instante ale
    demo-ului in acelasi timp.
\end{description}

\cmdchunk{../launch_realtime_demo.cmd}{10}{25}

\begin{description}
    \item[Liniile 10--17.] Daca lock-ul exista, scriptul verifica prin PowerShell
    daca mai ruleaza un proces care contine \texttt{src.tools.realtime\_demo} in
    linia de comanda. Daca procesul nu mai exista, lock-ul vechi este sters.
    \item[Liniile 20--25.] Scriptul incearca sa creeze lock-ul. Daca nu poate,
    concluzioneaza ca demo-ul este deja pornit si se opreste elegant, cu mesaj
    pentru utilizator.
\end{description}

\cmdchunk{../launch_realtime_demo.cmd}{27}{43}

\begin{description}
    \item[Liniile 27--34.] Verifica existenta checkpoint-ului. Interfata poate
    porni doar in modul MaskNet daca fisierul \path{masknet_best.pth} exista.
    \item[Liniile 36--40.] Selecteaza executabilul Python. Initial foloseste
    \texttt{python}, apoi prefera mediul virtual local daca exista
    \path{.venv/Scripts/python.exe} si acesta raspunde la \texttt{--version}.
    \item[Liniile 42--43.] Creeaza directoarele in care se salveaza recording-uri
    si snapshot-uri.
\end{description}

\cmdchunk{../launch_realtime_demo.cmd}{45}{62}

\begin{description}
    \item[Liniile 45--50.] Afiseaza in consola informatii despre proiect si
    interpretul Python ales, astfel incat utilizatorul sa poata diagnostica
    usor mediul de rulare.
    \item[Liniile 52--59.] Porneste modulul Python \texttt{src.tools.realtime\_demo}
    cu metoda \texttt{masknet}, checkpoint explicit, interfata vizuala, interval
    de statistici de 2 secunde si directoare dedicate pentru rezultate.
    \item[Liniile 61--62.] La final, lock-ul este sters, se afiseaza un mesaj de
    oprire si consola ramane deschisa prin \texttt{pause}.
\end{description}

\chapter{Entrypoint-ul Python}

\filetitle{src/tools/realtime_demo.py}

\chunk{../src/tools/realtime_demo.py}{1}{25}

\begin{description}
    \item[Liniile 1--3.] Docstring-ul explica scopul fisierului: acesta este
    punctul de intrare pentru demo-ul microfon-difuzoare.
    \item[Liniile 5--8.] Activeaza adnotarile amanate, importa \texttt{argparse}
    pentru CLI si \texttt{Path} pentru cai de fisiere.
    \item[Liniile 10--24.] Importa PyTorch, configuratia proiectului, clasele
    real-time si logger-ul. Observam ca entrypoint-ul nu implementeaza direct
    procesarea, ci doar compune componente deja definite in pachetul
    \texttt{src.tools.realtime}.
\end{description}

\chunk{../src/tools/realtime_demo.py}{27}{87}

\begin{description}
    \item[Liniile 27--30.] Defineste functia \texttt{parse\_args}, care construieste
    parserul de argumente pentru aplicatia real-time.
    \item[Liniile 31--41.] Argumentul \texttt{--method} limiteaza metodele la
    variantele acceptate: MaskNet, Spectral Subtraction, Wiener si bypass.
    Argumentul \texttt{--checkpoint} este optional, dar devine necesar daca se
    foloseste MaskNet fara checkpoint implicit disponibil.
    \item[Liniile 43--48.] Definesc dispozitivul PyTorch, sample rate-ul,
    dimensiunea blocului si contextul istoric. Valorile sunt exprimate in
    milisecunde pentru ergonomie, apoi sunt convertite mai tarziu in esantioane.
    \item[Liniile 52--60.] Definesc coada algoritmica, raportul dry/wet,
    dimensiunea cozilor interne si latenta ceruta bibliotecii \texttt{sounddevice}.
    \item[Liniile 61--66.] Controleaza raportarea statisticilor, dispozitivele
    audio si modul de vizualizare/recording.
    \item[Liniile 67--85.] Permit alegerea directoarelor pentru recording si
    snapshot-uri, durata istoricului vizual si modul fullscreen.
    \item[Linia 86.] Intoarce obiectul \texttt{Namespace}; de aici incolo,
    functia \texttt{main} foloseste toate optiunile intr-un mod centralizat.
\end{description}

\chunk{../src/tools/realtime_demo.py}{89}{167}

\begin{description}
    \item[Liniile 89--96.] Pornesc functia principala si verifica dependinta
    \texttt{sounddevice}. Daca biblioteca lipseste, aplicatia opreste rularea cu
    un mesaj clar.
    \item[Liniile 98--100.] Modul \texttt{--list-devices} afiseaza dispozitivele
    audio si iese fara a porni stream-ul.
    \item[Liniile 102--106.] Pentru MaskNet, alege checkpoint-ul implicit sau pe
    cel furnizat de utilizator si verifica existenta lui.
    \item[Liniile 108--113.] Alege dispozitivul PyTorch si converteste
    \texttt{block-ms}, \texttt{context-ms} si \texttt{tail-ms} in esantioane.
    \item[Liniile 115--124.] Creeaza obiectul \texttt{StreamingEnhancer}, adica
    stratul care va aplica metoda de denoising asupra blocurilor audio.
    \item[Linia 125.] Creeaza \texttt{FileDemoSource}, sursa alternativa pentru
    demonstratii sintetice \texttt{speech+noise}.
    \item[Liniile 127--135.] Selecteaza perechea de dispozitive audio si construieste
    liste scurte de candidati pentru combobox-urile din UI.
    \item[Liniile 137--139.] Daca utilizatorul cere recording, se creeaza sesiunea
    de inregistrare si se porneste colectarea blocurilor audio.
    \item[Liniile 141--163.] Daca este activ \texttt{--visualize}, se creeaza
    dashboard-ul Tkinter. In acest moment UI-ul primeste sample rate-ul, metoda
    initiala, checkpoint-ul, optiunile audio si starea sursei sintetice.
    \item[Liniile 165--166.] Se initializeaza etichetele de fisiere si dispozitivele
    active in UI.
    \item[Liniile 168--187.] Se afiseaza o nota practica despre casti, apoi se
    creeaza si ruleaza \texttt{RealtimeDenoiseApp}. Acesta este momentul in care
    controlul trece de la configurare la bucla real-time.
\end{description}

\chapter{Pachetul realtime}

\filetitle{src/tools/realtime/__init__.py}

\chunk{../src/tools/realtime/__init__.py}{1}{30}

\begin{description}
    \item[Linia 1.] Docstring-ul marcheaza pachetul ca un set de module
    compozabile pentru demo-ul real-time.
    \item[Liniile 3--14.] Reexporta utilitare din \texttt{common.py}: sursa de
    fisiere, recording-ul, statisticile, selectia dispozitivelor si accesul la
    \texttt{sounddevice}.
    \item[Linia 15.] Reexporta motorul audio si enhancer-ul.
    \item[Linia 16.] Reexporta interfata grafica.
    \item[Liniile 18--30.] Defineste \texttt{\_\_all\_\_}, adica lista simbolurilor
    considerate API public al pachetului. Aceasta lista face importurile din
    \texttt{realtime\_demo.py} mai curate si mai stabile.
\end{description}

\chapter{Utilitare comune}

\filetitle{src/tools/realtime/common.py}

\chunk{../src/tools/realtime/common.py}{1}{28}

\begin{description}
    \item[Liniile 1--3.] Docstring-ul defineste rolul fisierului: functii comune
    si utilitare pentru dispozitivele audio.
    \item[Liniile 5--15.] Importa module pentru logging, fire de executie,
    dataclass-uri, timp, cai de fisiere, NumPy, SoundFile si resampling.
    \item[Liniile 17--18.] Importa configuratia proiectului si pipeline-ul de
    augmentare, folosit de sursa sintetica \texttt{speech+noise}.
    \item[Liniile 20--23.] Incearca sa importe \texttt{sounddevice}. Daca biblioteca
    lipseste, variabila \texttt{sd} devine \texttt{None}; aplicatia poate raporta
    problema mai tarziu fara sa cada la import.
    \item[Liniile 26--28.] Definesc logger-ul si pragurile pentru detectarea
    microfonului aproape silentios.
\end{description}

\chunk{../src/tools/realtime/common.py}{30}{76}

\begin{description}
    \item[Liniile 30--48.] \texttt{resolve\_default\_checkpoint} cauta un checkpoint
    local potrivit. Prioritizeaza experimentele recomandate, apoi orice fisier
    \path{masknet_best.pth}. Daca nu gaseste nimic, ridica o eroare clara.
    \item[Liniile 50--51.] \texttt{ms\_to\_samples} transforma durata in milisecunde
    in numar de esantioane, folosind sample rate-ul.
    \item[Liniile 54--61.] \texttt{normalize\_audio\_block} curata valorile NaN/Inf,
    limiteaza varful la \texttt{0.98} si taie blocul in intervalul \([-1,1]\).
    \item[Liniile 63--69.] \texttt{rms\_dbfs} calculeaza nivelul RMS in dBFS pentru
    UI. Pragul minim evita logaritmul din zero.
    \item[Liniile 71--76.] \texttt{block\_rms} intoarce RMS liniar, folosit pentru
    euristica de microfon silentios.
\end{description}

\chunk{../src/tools/realtime/common.py}{78}{167}

\begin{description}
    \item[Liniile 78--101.] \texttt{choose\_audio\_file} deschide un dialog Tkinter
    pentru selectarea fisierelor audio. Functia este izolata aici pentru ca
    motorul sa poata cere fisiere fara a cunoaste detaliile UI.
    \item[Liniile 105--126.] \texttt{StreamStats} este un \texttt{dataclass} care
    acumuleaza numarul de callback-uri, blocuri procesate, underflow-uri,
    overflows si timpi de procesare. Proprietatea \texttt{mean\_processing\_time}
    evita impartirea la zero.
    \item[Liniile 129--167.] \texttt{RecordingSession} colecteaza blocuri raw si
    enhanced. Metodele \texttt{start}, \texttt{stop} si \texttt{toggle} controleaza
    starea, \texttt{add} copiaza blocuri doar cand recording-ul este activ, iar
    \texttt{save} concateneaza blocurile si scrie doua fisiere WAV.
\end{description}

\chunk{../src/tools/realtime/common.py}{170}{231}

\begin{description}
    \item[Liniile 170--214.] \texttt{FileDemoSource} defineste sursa alternativa
    \texttt{speech+noise}. Preset-urile \texttt{neutral}, \texttt{office},
    \texttt{street} si \texttt{hard} stabilesc SNR-ul si distorsiunile folosite
    in demonstratii.
    \item[Liniile 216--231.] Constructorul initializeaza sample rate-ul,
    augmenter-ul, fisierele incarcate, pozitia de playback, modul sursei,
    distorsiunile si un \texttt{Lock}. Lock-ul protejeaza starea sursei cand UI-ul
    si callback-ul audio o acceseaza simultan.
\end{description}

\chunk{../src/tools/realtime/common.py}{232}{317}

\begin{description}
    \item[Liniile 232--246.] \texttt{load\_speech\_file} si \texttt{load\_noise\_file}
    citesc fisierele, le normalizeaza, reseteaza pozitia si reconstruiesc mixul.
    \item[Liniile 248--260.] Controleaza modul sursei si play/pause-ul pentru
    fisiere. Toate modificarile sunt protejate de lock.
    \item[Liniile 262--280.] Actualizeaza distorsiunile, SNR-ul si preset-ul de
    scenariu. Orice schimbare reconstruieste mixul, pentru ca urmatoarele blocuri
    livrate callback-ului sa reflecte noua stare.
    \item[Liniile 283--303.] \texttt{get\_block} livreaza blocuri audio din mixul
    sintetic. Daca ajunge la finalul fisierului, continua circular prin concatenare
    cu inceputul semnalului.
    \item[Liniile 305--317.] Expun starea sursei: daca este pregatita si ce fisiere
    sunt incarcate. UI-ul foloseste aceste informatii pentru etichete si status.
\end{description}

\chunk{../src/tools/realtime/common.py}{319}{386}

\begin{description}
    \item[Liniile 319--325.] \texttt{\_load\_audio} citeste fisierul, converteste
    stereo la mono, resampleaza la sample rate-ul demo-ului si normalizeaza.
    \item[Liniile 327--331.] \texttt{\_resample\_audio} foloseste \texttt{resample\_poly}
    cu raport redus prin cel mai mare divizor comun, ceea ce evita resampling-ul
    inutil de scump.
    \item[Liniile 333--353.] \texttt{\_apply\_selected\_distortions} parcurge
    distorsiunile in ordine fixa si aplica doar cele activate.
    \item[Liniile 355--361.] \texttt{\_match\_length} potriveste lungimea zgomotului
    cu lungimea vorbirii, fie prin taiere, fie prin repetare.
    \item[Liniile 363--371.] \texttt{\_mix\_at\_snr} calculeaza factorul de scalare
    al zgomotului astfel incat mixul sa respecte SNR-ul dorit.
    \item[Liniile 373--386.] \texttt{\_rebuild\_mix} aplica distorsiunile pe vorbire,
    potriveste zgomotul si construieste semnalul final care va inlocui microfonul
    cand sursa este \texttt{speech+noise}.
\end{description}

\chunk{../src/tools/realtime/common.py}{388}{558}

\begin{description}
    \item[Liniile 388--402.] \texttt{print\_devices} listeaza toate dispozitivele
    audio disponibile, cu numarul de canale si sample rate-ul implicit.
    \item[Liniile 405--414.] \texttt{\_device\_keyword\_bonus} acorda scoruri pozitive
    sau negative pe baza numelui dispozitivului. De exemplu, \texttt{microphone}
    ajuta la input, iar \texttt{virtual} penalizeaza candidatul.
    \item[Liniile 416--441.] \texttt{\_score\_device\_candidate} combina scorul pe
    cuvinte-cheie cu sample rate-ul, dispozitivul implicit si mici penalizari
    pentru host API-uri mai putin potrivite.
    \item[Liniile 443--461.] \texttt{pick\_best\_device} alege cel mai bun input sau
    output disponibil, folosind scorurile de mai sus.
    \item[Liniile 463--469.] \texttt{resolve\_device} interpreteaza cererea
    utilizatorului: ID numeric, nume exact sau selectie automata.
    \item[Liniile 472--529.] \texttt{pick\_best\_device\_pair} incearca sa aleaga
    o pereche input-output compatibila. Daca poate, prefera dispozitive din acelasi
    host API, deoarece acestea sunt mai stabile in stream-uri simultane.
    \item[Liniile 531--540.] \texttt{describe\_device} transforma ID-ul intr-un text
    usor de afisat in loguri si UI.
    \item[Liniile 543--558.] \texttt{list\_device\_candidates} intoarce primii
    candidati pentru combobox-urile de input/output.
\end{description}

\chapter{Motorul real-time}

\filetitle{src/tools/realtime/engine.py}

\chunk{../src/tools/realtime/engine.py}{1}{30}

\begin{description}
    \item[Liniile 1--3.] Docstring-ul fixeaza rolul fisierului: motorul audio
    real-time.
    \item[Liniile 5--15.] Importa module pentru cozi, fire, timp, cai, NumPy,
    PyTorch si metodele de enhancement.
    \item[Liniile 16--27.] Importa utilitarele comune, inclusiv logger-ul,
    recorderul, sursa de fisiere, normalizarea si \texttt{sounddevice}.
    \item[Linia 29.] Importa clasa UI, dar motorul o foloseste optional; poate
    rula si fara interfata vizuala.
\end{description}

\chunk{../src/tools/realtime/engine.py}{32}{122}

\begin{description}
    \item[Liniile 32--56.] \texttt{StreamingEnhancer} pastreaza metoda curenta,
    dispozitivul PyTorch, dimensiunea blocului, contextul, taierea de coada,
    raportul dry/wet si checkpoint-ul. Bufferul \texttt{input\_buffer} este o
    fereastra glisanta folosita pentru procesare stabila.
    \item[Liniile 58--64.] \texttt{set\_method} permite schimbarea metodei live.
    Daca metoda devine \texttt{masknet}, modelul este incarcat lazy doar la prima
    utilizare.
    \item[Liniile 66--75.] Seteaza modul VAD, raportul dry/wet si ruleaza un
    warmup cu liniste pentru reducerea latentei primului bloc.
    \item[Liniile 77--102.] \texttt{process\_block} normalizeaza forma blocului,
    actualizeaza bufferul glisant, proceseaza fereastra completa, extrage blocul
    de iesire din zona stabila si aplica mixul dry/wet.
    \item[Liniile 104--122.] \texttt{\_enhance\_window} alege backend-ul efectiv:
    MaskNet, Spectral Subtraction, Wiener sau bypass. La final aliniaza lungimea
    rezultatului si normalizeaza blocul.
\end{description}

\chunk{../src/tools/realtime/engine.py}{125}{165}

\begin{description}
    \item[Liniile 125--157.] Constructorul \texttt{RealtimeDenoiseApp} primeste toate
    componentele deja create: enhancer-ul, sample rate-ul, dispozitivele, UI-ul,
    recorderul, sursa de fisiere si gain-ul de iesire.
    \item[Liniile 159--165.] Creeaza statisticile, evenimentul de oprire, doua cozi
    bounded, firul worker si starea folosita pentru detectarea microfonului
    silentios. Cozile limitate sunt importante pentru a evita acumularea latentei.
\end{description}

\chunk{../src/tools/realtime/engine.py}{167}{239}

\begin{description}
    \item[Liniile 167--181.] \texttt{run} face warmup, porneste firul worker si scrie
    informatii de log despre metoda, sample rate, block size, context si dispozitive.
    \item[Liniile 182--196.] Pregateste bucla principala si retine ultima pereche
    functionala de dispozitive, pentru fallback.
    \item[Liniile 197--215.] Deschide separat \texttt{InputStream} si \texttt{OutputStream}.
    Fiecare stream are callback propriu, ceea ce separa clar citirea de redare.
    \item[Liniile 216--239.] In interiorul stream-urilor, bucla doarme scurt si apoi
    verifica periodic comenzile venite din UI.
\end{description}

\chunk{../src/tools/realtime/engine.py}{240}{369}

\begin{description}
    \item[Liniile 240--281.] Consuma cererile UI: schimbare metoda, schimbare VAD,
    schimbare sursa, scenariu, SNR, dry/wet, gain si distorsiuni.
    \item[Liniile 282--310.] Proceseaza play/pause pentru fisiere, schimbarea
    dispozitivelor, cererile de fisiere, recording-ul si snapshot-ul.
    \item[Liniile 311--324.] Reimprospateaza UI-ul si opreste demo-ul daca fereastra
    a fost inchisa sau utilizatorul a cerut stop.
    \item[Liniile 325--334.] Scrie periodic statistici in log.
    \item[Liniile 335--369.] Trateaza erorile de stream audio. Daca noua pereche de
    dispozitive esueaza, revine la ultima pereche valida; altfel ramane intr-o
    bucla de asteptare in care UI-ul poate continua sa ceara o alta pereche.
\end{description}

\chunk{../src/tools/realtime/engine.py}{372}{447}

\begin{description}
    \item[Liniile 372--400.] \texttt{\_process\_file\_requests} deschide dialoguri
    pentru fisiere speech/noise, incarca sursa sintetica, porneste playback-ul si
    sincronizeaza etichetele din UI.
    \item[Liniile 402--447.] \texttt{\_audio\_callback} este o versiune combinata
    input-output, pastrata pentru compatibilitate conceptuala. Ea arata modelul
    initial full-duplex: citeste input, pune in coada, extrage output, aplica gain,
    trimite catre recorder si UI.
\end{description}

\chunk{../src/tools/realtime/engine.py}{449}{510}

\begin{description}
    \item[Liniile 449--476.] \texttt{\_input\_callback} este callback-ul actual de
    intrare. Copiaza blocul din microfon sau din sursa sintetica, actualizeaza
    ultimul bloc vazut, verifica linistea si il introduce in \texttt{input\_queue}.
    Daca coada este plina, arunca un bloc vechi pentru a mentine latenta mica.
    \item[Liniile 478--510.] \texttt{\_output\_callback} extrage blocul procesat din
    \texttt{output\_queue}. Daca nu exista bloc, reda liniste. Apoi aplica gain-ul,
    scrie in bufferul de iesire si trimite perechea raw/enhanced catre recorder si
    UI.
\end{description}

\chunk{../src/tools/realtime/engine.py}{512}{578}

\begin{description}
    \item[Liniile 512--536.] \texttt{\_update\_input\_activity\_warning} detecteaza
    un microfon aproape mut. Avertismentul se aplica doar cand sursa activa este
    microfonul, nu fisierul sintetic.
    \item[Liniile 538--560.] \texttt{\_processor\_loop} este firul worker. Preia
    blocuri din \texttt{input\_queue}, le proceseaza prin \texttt{StreamingEnhancer},
    masoara timpul si pune rezultatul in \texttt{output\_queue}.
    \item[Liniile 562--578.] \texttt{\_log\_stats} transforma timpii in milisecunde
    si scrie un rezumat: callback-uri, blocuri procesate, timpi medii/maximi,
    drop-uri, underflow-uri si avertismente.
\end{description}

\chapter{Interfata grafica}

\filetitle{src/tools/realtime/ui.py}

\chunk{../src/tools/realtime/ui.py}{1}{16}

\begin{description}
    \item[Liniile 1--3.] Docstring-ul spune ca fisierul implementeaza dashboard-ul
    Tkinter pentru comparatia raw vs enhanced.
    \item[Liniile 5--13.] Importa fire, timp, timestamp-uri, cai, NumPy,
    configuratia proiectului, VAD-ul si functiile de afisare a dispozitivelor si
    nivelurilor RMS.
    \item[Liniile 16--17.] Declara clasa \texttt{LiveComparisonUI}, punctul central
    al interfetei.
\end{description}

\chunk{../src/tools/realtime/ui.py}{17}{99}

\begin{description}
    \item[Liniile 17--38.] Semnatura constructorului primeste configuratia initiala:
    sample rate, durata bufferului afisat, metoda, recording, checkpoint,
    fullscreen, VAD, sursa, distorsiuni, SNR, dry/wet, gain, scenariu si listele
    de dispozitive.
    \item[Liniile 39--56.] Initializeaza bufferele audio, nivelurile RMS, metoda
    curenta, metoda pending, VAD-ul, sursa, scenariul, starea de stop, recording-ul
    si snapshot-ul.
    \item[Liniile 57--70.] Defineste lista de distorsiuni disponibile si creeaza
    hartile curente/pending. Acest mecanism permite UI-ului sa modifice distorsiuni
    fara sa atinga direct motorul.
    \item[Liniile 71--88.] Initializeaza cererile pentru fisiere, starea de playback,
    SNR-ul, dry/wet-ul, gain-ul si etichetele fisierelor incarcate.
    \item[Liniile 89--99.] Pastreaza optiunile de dispozitive, starea activa,
    mesajul de status si lock-ul care protejeaza bufferele audio afisate.
\end{description}

\chunk{../src/tools/realtime/ui.py}{101}{139}

\begin{description}
    \item[Liniile 101--104.] Importa Tkinter, \texttt{ttk} si Matplotlib local, in
    constructor. Astfel, modulul poate fi importat fara a deschide implicit UI-ul.
    \item[Liniile 106--119.] Creeaza fereastra principala, seteaza titlul,
    dimensiunea, modul fullscreen si handlerul de inchidere.
    \item[Liniile 121--139.] Configureaza stilurile \texttt{ttk}: panouri, etichete,
    taburi, butoane si bare de progres.
\end{description}

\chunk{../src/tools/realtime/ui.py}{141}{214}

\begin{description}
    \item[Liniile 141--151.] Creeaza layout-ul principal cu doua coloane: zona de
    monitorizare si zona de control.
    \item[Liniile 153--167.] Creeaza headerul si statusul de semnal, apoi panoul
    de sumar al sesiunii.
    \item[Liniile 169--190.] Creeaza figura Matplotlib principala cu doua waveform-uri
    si doua spectrograme.
    \item[Liniile 191--211.] Initializeaza imaginile spectrogramelor cu matrici zero,
    palete diferite si scala dB intre \(-80\) si \(0\).
    \item[Liniile 212--214.] Integreaza figura Matplotlib in Tkinter prin
    \texttt{FigureCanvasTkAgg}.
\end{description}

\chunk{../src/tools/realtime/ui.py}{216}{346}

\begin{description}
    \item[Liniile 216--235.] Creeaza figura VAD si panoul din dreapta al interfetei.
    \item[Liniile 238--251.] Creeaza zona de niveluri RMS pentru input si output,
    folosind bare de progres si valori dBFS.
    \item[Liniile 253--276.] Creeaza notebook-ul cu taburi: Audio, Processing,
    Files si Augment. Initializeaza variabilele Tkinter asociate controalelor.
    \item[Liniile 278--291.] Tabul Audio permite alegerea sursei si a dispozitivelor
    de intrare/iesire.
    \item[Liniile 293--307.] Tabul Processing permite alegerea metodei, VAD-ului,
    dry/wet-ului, gain-ului, snapshot-ului, recording-ului si opririi demo-ului.
    \item[Liniile 309--320.] Tabul Files permite incarcarea fisierelor de vorbire
    si zgomot si controlul play/pause.
    \item[Liniile 322--337.] Tabul Augment permite alegerea scenariului, SNR-ului
    si distorsiunilor.
    \item[Liniile 339--346.] Ataseaza \texttt{trace}-uri variabilelor Tkinter, astfel
    incat modificarile din combobox-uri sa seteze cereri pending.
\end{description}

\chunk{../src/tools/realtime/ui.py}{347}{429}

\begin{description}
    \item[Liniile 347--351.] \texttt{\_find\_device\_label} gaseste eticheta UI
    pentru un ID de dispozitiv sau intoarce prima optiune disponibila.
    \item[Liniile 353--358.] \texttt{\_on\_close} marcheaza UI-ul ca inchis si distruge
    fereastra.
    \item[Liniile 360--369.] Handlerele pentru stop, recording, snapshot si aplicarea
    dispozitivelor seteaza flag-uri pending.
    \item[Liniile 372--411.] Handlerele pentru metoda, sursa, input/output, VAD,
    scenariu, SNR, dry/wet si gain actualizeaza variabilele pending.
    \item[Liniile 415--421.] Handlerele pentru fisiere si play/pause seteaza cereri
    pe care motorul le va consuma.
    \item[Liniile 424--429.] Cand o distorsiune este bifata, UI-ul reconstruieste
    harta pending a tuturor distorsiunilor.
\end{description}

\chunk{../src/tools/realtime/ui.py}{430}{527}

\begin{description}
    \item[Liniile 430--437.] \texttt{push} primeste blocurile raw si enhanced din
    callback-ul de iesire. Actualizeaza bufferele circulare si nivelurile RMS sub
    protectia lock-ului.
    \item[Liniile 439--527.] Metodele \texttt{consume\_*} implementeaza modelul
    \texttt{pending-consume}. Fiecare metoda verifica daca exista o schimbare
    ceruta de UI, intoarce noua valoare si reseteaza flag-ul. Motorul apeleaza
    aceste metode intr-un punct sigur al buclei real-time.
\end{description}

\chunk{../src/tools/realtime/ui.py}{529}{603}

\begin{description}
    \item[Liniile 529--579.] Metodele \texttt{set\_*} sincronizeaza UI-ul cu starea
    aplicata de motor: recording, fisiere incarcate, play/pause, dispozitive,
    mesaj de stream, SNR, dry/wet, gain, scenariu si distorsiuni.
    \item[Liniile 581--603.] \texttt{\_compute\_display\_spectrogram} calculeaza
    spectrograma pentru afisare: pad daca semnalul este scurt, fereastra Hann,
    FFT real, conversie in dB si limitare intre \(-80\) si \(0\).
\end{description}

\chunk{../src/tools/realtime/ui.py}{605}{709}

\begin{description}
    \item[Liniile 605--616.] \texttt{refresh} copiaza bufferele sub lock si actualizeaza
    waveform-urile.
    \item[Liniile 618--630.] Recalculeaza spectrogramele si curba VAD soft, apoi
    actualizeaza axele si titlul panoului VAD.
    \item[Liniile 632--672.] Actualizeaza nivelurile RMS, sumarul sesiunii, statusul
    checkpoint-ului, sursa, VAD-ul, scenariul, SNR-ul, dry/wet-ul, gain-ul,
    latenta si dispozitivele.
    \item[Liniile 674--696.] Actualizeaza indicatorul de semnal, titlurile graficelor
    si redeseneaza Matplotlib la cel mult aproximativ 12 cadre pe secunda.
    \item[Liniile 698--704.] \texttt{save\_snapshot} salveaza figura principala ca
    PNG intr-un director de rezultate.
    \item[Liniile 706--709.] \texttt{close} distruge fereastra daca nu a fost deja
    inchisa.
\end{description}

\chapter{Concluzie}

Fișierele real-time sunt organizate pe trei niveluri. Launcher-ul si
\texttt{realtime\_demo.py} pornesc aplicatia si compun obiectele. \texttt{common.py}
ofera infrastructura partajata pentru audio, fisiere, recording si dispozitive.
\texttt{engine.py} este responsabil de constrangerile real-time: callback-uri,
cozi, thread worker si procesare pe blocuri. \texttt{ui.py} expune toate aceste
stari intr-un dashboard interactiv si foloseste modelul \texttt{pending-consume}
pentru a trimite comenzi catre motor fara a bloca audio callback-ul.

Prin aceasta separare, codul ramane mai usor de explicat si de testat: UI-ul nu
proceseaza direct semnalul, motorul nu construieste widget-uri, iar entrypoint-ul
nu contine logica DSP. Fiecare fisier are o responsabilitate clara in lantul
real-time de captare, procesare, redare si vizualizare.

\end{document}
